
Our work intersects parameter-efficient fine-tuning, multi-task learning, and neural architecture search. We position our contribution as the first systematic measurement of LoRA rank oracle gap on multi-domain benchmarks.

\subsubsection{Parameter-Efficient Fine-Tuning}

LoRA (Hu et al., 2021) introduced low-rank adaptation, injecting trainable rank-decomposition matrices into frozen transformer layers. The method achieves comparable performance to full fine-tuning while training only 0.1-1\% of parameters, with rank r controlling capacity through O(d·r) parameter scaling. Follow-up work includes AdaLoRA (Zhang et al., 2023), which adapts rank dynamically during training, and QLoRA (Dettmers et al., 2023), which combines LoRA with quantization.

All these methods treat rank as a global hyperparameter chosen once and applied uniformly. Practitioners typically fix rank at 8 or 16 based on validation performance on a single task or average performance across a task suite. Our work measures the oracle gap from per-task rank selection, demonstrating that heterogeneous task distributions create opportunities that fixed configurations cannot capture.

Adapter-based methods (Houlsby et al., 2019; Pfeiffer et al., 2020) inject small bottleneck layers into transformers. AdapterHub (Pfeiffer et al., 2020) provides a framework for sharing and composing adapters across tasks, demonstrating that different tasks benefit from different adapter configurations. While this work shows task-specific adapter effectiveness, it does not quantify the oracle gap from capacity selection.

Prefix tuning and prompt-based methods (Li and Liang, 2021; Lester et al., 2021) prepend trainable vectors to input sequences. These methods also face the fixed-configuration problem: prefix length and depth are chosen globally. Our oracle gap measurement methodology could extend to these methods, though we focus on LoRA.

\subsubsection{Multi-Task Learning and Task Heterogeneity}

Multi-task learning (Caruana, 1997; Ruder, 2017) has recognized that different tasks may benefit from different model capacities or architectures. Recent work on task-conditional computation includes Mixture-of-Experts (Shazeer et al., 2017; Fedus et al., 2022), which routes inputs to different expert subnetworks during training.

These methods route during training to learn shared representations with task-specific specialization, not at deployment to select from pre-trained adapter configurations. Our work focuses on the deployment scenario: given a library of pre-trained adapters with different ranks, how much performance is lost by using a single fixed rank versus selecting per-task?

Task complexity analysis (Bingel and Søgaard, 2017; Talmor and Berant, 2019) has studied why some tasks are harder than others. Our finding that oracle rank selections vary across tasks aligns with this literature, extending it to the adapter capacity dimension.

\subsubsection{Neural Architecture Search and Hyperparameter Optimization}

Neural Architecture Search (Zoph and Le, 2017; Pham et al., 2018; Liu et al., 2019) automatically discovers optimal architectures for each task. AutoML systems (Feurer et al., 2015; Thornton et al., 2013) similarly optimize hyperparameters per task through extensive search.

These approaches validate that per-task optimization can improve over fixed configurations, implicitly measuring oracle-like gaps when comparing per-task tuned systems to global defaults. Our work extends this concept to the adapter rank dimension with a lightweight discrete configuration space (4 ranks).

Once-for-all networks (Cai et al., 2020) train a single super-network that supports multiple sub-architectures, enabling deployment-time selection without retraining. This paradigm aligns with training multiple adapter ranks once, then selecting at deployment based on task characteristics. However, once-for-all networks focus on edge device constraints, while we focus on multi-task performance optimization.

\subsubsection{Multi-Domain Benchmarks}

GLUE (Wang et al., 2018) and SuperGLUE (Wang et al., 2019) established multi-task evaluation for natural language understanding. XTREME (Hu et al., 2020) extends multi-domain evaluation to cross-lingual settings, covering 40 languages across 9 task types.

These benchmarks were designed for task diversity—measuring model robustness across linguistic phenomena. Our work leverages this diversity to test whether optimal adapter configurations vary across tasks. The distribution of oracle selections (5/4/4/4 across ranks) confirms that these benchmarks exhibit heterogeneity in optimal adapter capacity.

\subsubsection{Positioning Our Contribution}

Prior work has established that parameter-efficient fine-tuning methods work with fixed rank configurations, that different tasks have different characteristics, and that per-task optimization can improve performance at high computational cost.

Our contribution is the first systematic measurement of the oracle gap for LoRA adaptation on multi-domain NLP benchmarks. By training all rank-task configurations and measuring the oracle gap (15.09\% relative improvement, 11.62 percentage points absolute), we quantify what fixed-configuration approaches leave on the table. This measurement establishes an upper bound for future task-aware routing research.
